% answers.tex —— 2014 年·天津卷·理科 · 参考答案与详细解析片段
% 由 tools/generate_projects.py 依来源总账自 2014/2014_数学真题_天津理_试题.md + 2014/2014_数学真题_天津理_解析正文.md + 2014/2014_数学真题_天津理_答案要点.md 生成，勿手改（改源 MD 后重跑）


\begin{answer}
\gksolhead{1}{}

\textbf{题目}：

$i$ 是虚数单位，复数 $\dfrac{7+\mathrm i}{3+4\mathrm i}=$\mbox{（　　）}

（A）$1-\mathrm i$；（B）$-1+\mathrm i$；（C）$\dfrac{17}{25}+\dfrac{31}{25}\mathrm i$；（D）$-\dfrac{17}{7}+\dfrac{25}{7}\mathrm i$。

\Needspace{4\baselineskip}
\textbf{解析}

\Needspace{5\baselineskip}
分子、分母同乘分母的共轭复数 $3-4\mathrm i$：
\nopagebreak[4]
\[\begin{aligned}
\frac{7+\mathrm i}{3+4\mathrm i}
&=\frac{(7+\mathrm i)(3-4\mathrm i)}{(3+4\mathrm i)(3-4\mathrm i)}\\
&=\frac{21-28\mathrm i+3\mathrm i-4\mathrm i^2}{3^2+4^2}\\
&=\frac{25-25\mathrm i}{25}\\
&=1-\mathrm i .
\end{aligned}\]

选 A。

\gkanswer{答案：A。分子分母同乘 $3-4\mathrm i$，得 $\dfrac{25-25\mathrm i}{25}=1-\mathrm i$。}

\end{answer}
\begin{answer}
\gksolhead{2}{}

\textbf{题目}：

\Needspace{5\baselineskip}
设变量 $x$，$y$ 满足约束条件
\nopagebreak[4]
\[\begin{cases}
x+y-2\geqslant 0,\\
x-y-2\leqslant 0,\\
y\geqslant 1,
\end{cases}\]
则目标函数 $z=x+2y$ 的最小值为\mbox{（　　）}

（A）$2$；（B）$3$；（C）$4$；（D）$5$。

\Needspace{4\baselineskip}
\textbf{解析}

\Needspace{5\baselineskip}
由 $x+y-2\geqslant 0$ 得 $x\geqslant 2-y$，由 $x-y-2\leqslant 0$ 得 $x\leqslant 2+y$。可行域内的点同时满足
\nopagebreak[4]
\[2-y\leqslant x\leqslant 2+y,\qquad y\geqslant 1,\]
\Needspace{5\baselineskip}
故
\nopagebreak[4]
\[z=x+2y\geqslant (2-y)+2y=2+y\geqslant 3.\]
当 $x=1$，$y=1$ 时，$x+y=2$、$x-y=0$ 均满足约束，且 $z=1+2=3$。所以 $z$ 的最小值为 $3$。

选 B。

\gkanswer{答案：B。由 $x\geqslant 2-y$ 得 $z=x+2y\geqslant 2+y\geqslant 3$，在点 $(1,1)$ 处取到 $3$。}

\end{answer}
\begin{answer}
\gksolhead{3}{}

\textbf{题目}：

阅读下面的程序框图，运行相应的程序，输出的 $S$ 的值为\mbox{（　　）}
\nopagebreak[4]
\begin{center}% ============================================================
% f2014_tianjin_science_q03_flowchart.tex
% 2014 天津卷（理科）第 3 题：程序框图
% 来源 PDF：普通高考/2014/2014天津理.pdf
% 原页：第 1 页（A3），印面 923，第 3 题题干下方
% S=1，i=1；循环体 T=2i+1、S=S×T、i=i+1；i≥4 成立则输出 S，否则返回循环体。
% ============================================================
\begin{tikzpicture}[>=Stealth, line width=0.6pt, font=\small, scale=0.85,
  line cap=round, line join=round,
  term/.style={draw, rectangle, rounded corners=2.2pt, minimum width=1.9cm,
    minimum height=0.72cm, inner sep=2pt},
  proc/.style={draw, rectangle, minimum width=4.2cm, minimum height=0.72cm,
    inner sep=2pt},
  io/.style={draw, trapezium, trapezium left angle=68, trapezium right angle=112,
    minimum width=3.4cm, minimum height=0.78cm, inner sep=2pt},
  dec/.style={draw, diamond, aspect=2.1, minimum width=4.4cm, minimum height=1.4cm,
    inner sep=1pt}]
  \node[term] (start) at (0,0) {开始};
  \node[proc] (init) at (0,-1.45) {$S=1$，$i=1$};
  \node[proc] (body) at (0,-2.95) {$T=2i+1$};
  \node[proc] (mul)  at (0,-4.45) {$S=S\times T$};
  \node[proc] (inc)  at (0,-5.95) {$i=i+1$};
  \node[dec]  (test) at (0,-7.80) {$i\geqslant 4$？};
  \node[io]   (out)  at (0,-9.55) {输出 $S$};
  \node[term] (end)  at (0,-10.85) {结束};

  \draw[->] (start) -- (init);
  \draw[->] (init) -- (body);
  \draw[->] (body) -- (mul);
  \draw[->] (mul) -- (inc);
  \draw[->] (inc) -- (test);
  \draw[->] (test) -- (out);
  \draw[->] (out) -- (end);

  % “是”标注：判断框下沿与输出框之间，箭头右侧
  \node[right, inner sep=2pt] at (0.32,-8.78) {是};

  % “否”：从判断框左顶点引出，左侧上升后折向循环体顶边
  \draw[->] (test.west) -- (-3.25,-7.80) -- (-3.25,-2.30) -- (-0.95,-2.30)
    -- (-0.95,-2.59);
  \node[above, inner sep=2pt] at (-2.75,-7.80) {否};
\end{tikzpicture}
\end{center}
\nopagebreak[4]
（A）$15$；（B）$105$；（C）$245$；（D）$945$。

\Needspace{4\baselineskip}
\textbf{解析}

按框图逐步执行：

\begin{itemize}
\item $i=1$ 时，$T=2\times1+1=3$，$S=1\times3=3$，$i$ 变为 $2$；
\item $i=2$ 时，$T=2\times2+1=5$，$S=3\times5=15$，$i$ 变为 $3$；
\item $i=3$ 时，$T=2\times3+1=7$，$S=15\times7=105$，$i$ 变为 $4$。
\end{itemize}

此时 $i=4\geqslant 4$，判断成立，输出 $S=105$。

选 B。

\gkanswer{答案：B。依次算得 $S=3$（$i=1$）、$S=15$（$i=2$）、$S=105$（$i=3$），此后 $i=4$ 满足 $i\geqslant 4$，输出 $105$。}

\end{answer}
\begin{answer}
\gksolhead{4}{}

\textbf{题目}：

函数 $f(x)=\log_{\frac12}(x^2-4)$ 的单调递增区间是\mbox{（　　）}

（A）$(0,+\infty)$；（B）$(-\infty,0)$；（C）$(2,+\infty)$；（D）$(-\infty,-2)$。

\Needspace{4\baselineskip}
\textbf{解析}

\Needspace{5\baselineskip}
由 $x^2-4>0$ 得定义域
\nopagebreak[4]
\[x<-2\quad\text{\text{或}}\quad x>2.\]
令 $u=x^2-4$，则 $y=\log_{\frac12}u$。因为底数 $\dfrac12\in(0,1)$，$y$ 随 $u$ 增大而减小；而 $u=x^2-4$ 在 $(-\infty,-2)$ 上随 $x$ 增大而减小。两个“减小”复合为“增大”，故 $f(x)$ 在 $(-\infty,-2)$ 上单调递增。在 $(2,+\infty)$ 上 $u$ 递增，$f(x)$ 递减。

因此单调递增区间是 $(-\infty,-2)$。

选 D。

\gkanswer{答案：D。定义域为 $|x|>2$；内层 $x^2-4$ 在 $(-\infty,-2)$ 上递减，外层 $\log_{\frac12}$ 递减，复合后递增。}

\end{answer}
\begin{answer}
\gksolhead{5}{}

\textbf{题目}：

已知双曲线 $\dfrac{x^2}{a^2}-\dfrac{y^2}{b^2}=1$（$a>0$，$b>0$）的一条渐近线平行于直线 $l:y=2x+10$，双曲线的一个焦点在直线 $l$ 上，则双曲线的方程为\mbox{（　　）}

（A）$\dfrac{x^2}{5}-\dfrac{y^2}{20}=1$；（B）$\dfrac{x^2}{20}-\dfrac{y^2}{5}=1$；（C）$\dfrac{3x^2}{25}-\dfrac{3y^2}{100}=1$；（D）$\dfrac{3x^2}{100}-\dfrac{3y^2}{25}=1$。

\Needspace{4\baselineskip}
\textbf{解析}

\Needspace{5\baselineskip}
双曲线的渐近线为 $y=\pm\dfrac ba x$。其中一条平行于 $y=2x+10$，故
\nopagebreak[4]
\[\frac ba=2,\qquad b=2a.\]
\Needspace{5\baselineskip}
焦点为 $(\pm c,0)$，在 $l:y=2x+10$ 上令 $y=0$ 得 $x=-5$，所以 $c=5$，且
\nopagebreak[4]
\[a^2+b^2=c^2=25.\]
\Needspace{5\baselineskip}
代入 $b=2a$：
\nopagebreak[4]
\[a^2+4a^2=25,\qquad a^2=5,\qquad b^2=20.\]
\Needspace{5\baselineskip}
故双曲线方程为
\nopagebreak[4]
\[\frac{x^2}{5}-\frac{y^2}{20}=1.\]

选 A。

\gkanswer{答案：A。由 $\dfrac ba=2$ 与焦点 $(-5,0)$ 得 $c=5$，$a^2=5$，$b^2=20$，即 $\dfrac{x^2}{5}-\dfrac{y^2}{20}=1$。}

\end{answer}
\begin{answer}
\gksolhead{6}{}

\textbf{题目}：

如图，$\triangle ABC$ 是圆的内接三角形，$\angle BAC$ 的平分线交圆于点 $D$，交 $BC$ 于 $E$，过点 $B$ 的圆的切线与 $AD$ 的延长线交于点 $F$。在上述条件下，给出下列四个结论：① $BD$ 平分 $\angle CBF$；② $FB^2=FD\cdot FA$；③ $AE\cdot CE=BE\cdot DE$；④ $AF\cdot BD=AB\cdot BF$。则所有正确结论的序号是\mbox{（　　）}
\nopagebreak[4]
\begin{center}% ============================================================
% f2014_tianjin_science_q06_circle.tex
% 2014 天津卷（理科）第 6 题：圆内接三角形与切线
% 来源 PDF：普通高考/2014/2014天津理.pdf
% 原页：第 1 页（A3），印面 923，第 6 题题干下方
% D 为 ∠BAC 平分线与圆的交点（弧 BC 中点），E=AD∩BC，
% F 为过 B 的切线与 AD 延长线的交点。取 A=70°、B=210°、C=330°、D=270°。
% ============================================================
\begin{tikzpicture}[>=Stealth, line width=0.7pt, font=\small,
  line cap=round, line join=round, scale=0.78]
  % 圆（半径 2）
  \draw[line width=0.7pt] (0,0) circle (2);
  \coordinate (A) at (0.684,1.879);
  \coordinate (B) at (-1.732,-1.000);
  \coordinate (C) at (1.732,-1.000);
  \coordinate (D) at (0,-2);
  \coordinate (E) at (0.176,-1.000);
  \coordinate (F) at (-0.270,-3.530);

  % 三角形与内接弦
  \draw (A) -- (B) -- (C) -- cycle;
  \draw (B) -- (D) -- (C);
  % 角平分线 AD 及延长线到 F
  \draw (A) -- (D);
  \draw (D) -- (F);
  % 过 B 的切线
  \draw (B) -- (F);

  \fill (A) circle (0.9pt) node[above right=-1pt] {$A$};
  \fill (B) circle (0.9pt) node[left=-1pt] {$B$};
  \fill (C) circle (0.9pt) node[right=-1pt] {$C$};
  \fill (D) circle (0.9pt) node[below right=-1pt] {$D$};
  \fill (E) circle (0.9pt) node[above left=-1pt] {$E$};
  \fill (F) circle (0.9pt) node[below=-1pt] {$F$};
\end{tikzpicture}
\end{center}
\nopagebreak[4]
（A）①②；（B）③④；（C）①②③；（D）①②④。

\Needspace{4\baselineskip}
\textbf{解析}

因为 $AD$ 平分 $\angle BAC$，所以 $\angle BAD=\angle DAC$，从而弧 $BD$ 与弧 $DC$ 相等，即 $D$ 是弧 $BC$（不含 $A$）的中点。

\Needspace{5\baselineskip}
① 弧 $BD$ 与弧 $DC$ 相等，而 $\angle CBD$ 与 $\angle BAD$ 分别对着弧 $CD$ 与弧 $BD$，由等弧上的圆周角相等得 $\angle CBD=\angle BAD$；又由弦切角定理，弦 $BD$ 与切线 $BF$ 的夹角 $\angle DBF$ 等于弧 $BD$ 所对的圆周角 $\angle BAD$。于是
\nopagebreak[4]
\[\angle CBD=\angle BAD=\angle DBF,\]
即 $BD$ 平分 $\angle CBF$。① 正确。

\Needspace{5\baselineskip}
② $F$ 在圆外，$FB$ 切圆于 $B$，$FAD$ 是过 $F$ 的割线，由切割线定理
\nopagebreak[4]
\[FB^2=FD\cdot FA.\]
② 正确。

\Needspace{5\baselineskip}
③ 点 $E$ 在圆内，由相交弦定理应为
\nopagebreak[4]
\[EA\cdot ED=EB\cdot EC,\]
而题中写成 $AE\cdot CE=BE\cdot DE$，两者一般不相等。③ 错误。

\Needspace{5\baselineskip}
④ 在 $\triangle ABF$ 与 $\triangle BDF$ 中，$\angle AFB=\angle BFD$（公共角）。又 $\angle BAF=\angle BAD$ 是弧 $BD$ 所对的圆周角，而 $\angle DBF$ 是弦 $BD$ 与切线的夹角，由弦切角定理 $\angle DBF=\angle BAD$。故
\nopagebreak[4]
\[\angle BAF=\angle DBF,\]
\Needspace{5\baselineskip}
从而 $\triangle ABF\backsim\triangle BDF$，对应边成比例：
\nopagebreak[4]
\[\frac{AB}{BD}=\frac{AF}{BF},\]
即 $AF\cdot BD=AB\cdot BF$。④ 正确。

所以正确的是 ①②④。

选 D。

\gkanswer{答案：D。弦切角定理与同弧圆周角给出 ①，切割线定理给出 ②，$\triangle ABF\backsim\triangle BDF$ 给出 ④；③ 应为 $AE\cdot DE=BE\cdot CE$。}

\end{answer}
\begin{answer}
\gksolhead{7}{}

\textbf{题目}：

设 $a$，$b\in\mathbf R$，则“$a>b$”是“$a|a|>b|b|$”的\mbox{（　　）}

（A）充分不必要条件；（B）必要不充分条件；（C）充要条件；（D）既不充分也不必要条件。

\Needspace{4\baselineskip}
\textbf{解析}

\Needspace{5\baselineskip}
设 $\varphi(t)=t|t|$，即
\nopagebreak[4]
\[\varphi(t)=\begin{cases}
t^2, & t\geqslant 0,\\
-t^2, & t<0.
\end{cases}\]
在 $[0,+\infty)$ 上 $\varphi(t)=t^2$ 严格递增；在 $(-\infty,0)$ 上 $\varphi(t)=-t^2$ 也严格递增，且两段在 $t=0$ 处连续、函数值同为 $0$。故 $\varphi$ 在 $\mathbf R$ 上严格递增。

\Needspace{5\baselineskip}
于是
\nopagebreak[4]
\[a>b\iff \varphi(a)>\varphi(b),\]
即“$a>b$”与“$a|a|>b|b|$”互为充要条件。

选 C。

\gkanswer{答案：C。$\varphi(t)=t|t|$ 在 $\mathbf R$ 上严格递增，故 $a>b$ 与 $\varphi(a)>\varphi(b)$ 等价。}

\end{answer}
\begin{answer}
\gksolhead{8}{}

\textbf{题目}：

已知菱形 $ABCD$ 的边长为 $2$，$\angle BAD=120^\circ$，点 $E$，$F$ 分别在边 $BC$，$DC$ 上，$\overrightarrow{BE}=\lambda\overrightarrow{BC}$，$\overrightarrow{DF}=\mu\overrightarrow{DC}$，若 $\overrightarrow{AE}\cdot\overrightarrow{AF}=1$，$\overrightarrow{CE}\cdot\overrightarrow{CF}=-\dfrac23$，则 $\lambda+\mu=$\mbox{（　　）}

（A）$\dfrac12$；（B）$\dfrac23$；（C）$\dfrac56$；（D）$\dfrac{7}{12}$。

\Needspace{4\baselineskip}
\textbf{解析}

\Needspace{5\baselineskip}
记 $\boldsymbol u=\overrightarrow{AB}$，$\boldsymbol v=\overrightarrow{AD}$。菱形边长为 $2$，$\angle BAD=120^\circ$，故
\nopagebreak[4]
\[|\boldsymbol u|=|\boldsymbol v|=2,\qquad \boldsymbol u\cdot\boldsymbol v=2\times2\times\cos120^\circ=-2.\]

\Needspace{5\baselineskip}
由 $\overrightarrow{BC}=\overrightarrow{AD}=\boldsymbol v$ 得
\nopagebreak[4]
\[\overrightarrow{AE}=\boldsymbol u+\lambda\boldsymbol v.\]
\Needspace{5\baselineskip}
由 $\overrightarrow{DC}=\overrightarrow{AB}=\boldsymbol u$ 得
\nopagebreak[4]
\[\overrightarrow{AF}=\boldsymbol v+\mu\boldsymbol u.\]
\Needspace{5\baselineskip}
于是
\nopagebreak[4]
\[\begin{aligned}
\overrightarrow{AE}\cdot\overrightarrow{AF}
&=(\boldsymbol u+\lambda\boldsymbol v)\cdot(\boldsymbol v+\mu\boldsymbol u)\\
&=\boldsymbol u\cdot\boldsymbol v+\mu|\boldsymbol u|^2+\lambda|\boldsymbol v|^2+\lambda\mu\,\boldsymbol v\cdot\boldsymbol u\\
&=-2+4\mu+4\lambda-2\lambda\mu.
\end{aligned}\]
\Needspace{5\baselineskip}
由 $\overrightarrow{AE}\cdot\overrightarrow{AF}=1$ 得
\nopagebreak[4]
\[4(\lambda+\mu)-2\lambda\mu=3. \tag{1}\]

\Needspace{5\baselineskip}
又 $\overrightarrow{AC}=\boldsymbol u+\boldsymbol v$，故
\nopagebreak[4]
\[\overrightarrow{CE}=\overrightarrow{AE}-\overrightarrow{AC}=(\lambda-1)\boldsymbol v,\qquad
\overrightarrow{CF}=\overrightarrow{AF}-\overrightarrow{AC}=(\mu-1)\boldsymbol u.\]
\Needspace{5\baselineskip}
由 $\overrightarrow{CE}\cdot\overrightarrow{CF}=-\dfrac23$ 得
\nopagebreak[4]
\[-2(\lambda-1)(\mu-1)=-\frac23,\]
\Needspace{5\baselineskip}
即
\nopagebreak[4]
\[(\lambda-1)(\mu-1)=\frac13,\qquad \lambda\mu-(\lambda+\mu)+1=\frac13,\]
\Needspace{5\baselineskip}
所以
\nopagebreak[4]
\[\lambda\mu=\lambda+\mu-\frac23. \tag{2}\]
\Needspace{5\baselineskip}
把（2）代入（1），记 $s=\lambda+\mu$：
\nopagebreak[4]
\[4s-2\left(s-\frac23\right)=3,\qquad 2s+\frac43=3,\qquad s=\frac56.\]
故 $\lambda+\mu=\dfrac56$。

选 C。

\gkanswer{答案：C。由 $\overrightarrow{AE}\cdot\overrightarrow{AF}=1$ 得 $4(\lambda+\mu)-2\lambda\mu=3$，由 $\overrightarrow{CE}\cdot\overrightarrow{CF}=-\dfrac23$ 得 $\lambda\mu=\lambda+\mu-\dfrac23$，解得 $\lambda+\mu=\dfrac56$。}

\end{answer}
\begin{answer}
\gksolhead{9}{}

\textbf{题目}：

某大学为了解在校本科生对参加某项社会实践活动的意向，拟采用分层抽样的方法，从该校四个年级的本科生中抽取一个容量为 $300$ 的样本进行调查。已知该校一年级、二年级、三年级、四年级的本科生人数之比为 $4:5:5:6$，则应从一年级本科生中抽取\gkblank{}名学生。

\Needspace{4\baselineskip}
\textbf{解析}

\Needspace{5\baselineskip}
四个年级的人数之比为 $4:5:5:6$，总份数为
\nopagebreak[4]
\[4+5+5+6=20.\]
\Needspace{5\baselineskip}
样本量为 $300$，故每份对应
\nopagebreak[4]
\[\frac{300}{20}=15\]
\Needspace{5\baselineskip}
人。一年级占 $4$ 份，应抽取
\nopagebreak[4]
\[4\times15=60\]
人。

填 $60$。

\gkanswer{答案：$60$。总份数 $20$，每份 $15$ 人，一年级 $4$ 份即 $60$ 人。}

\end{answer}
\begin{answer}
\gksolhead{10}{}

\textbf{题目}：

已知一个几何体的三视图如图所示（单位：m），则该几何体的体积为\gkblank{}$\mathrm m^3$。
\nopagebreak[4]
\begin{center}% ============================================================
% f2014_tianjin_science_q10_views.tex
% 2014 天津卷（理科）第 10 题：三视图
% 来源 PDF：普通高考/2014/2014天津理.pdf
% 原页：第 1 页（A3），印面 923，第 10 题题干下方
% 正视图、侧视图相同：宽 2、高 4 的矩形上加底宽 4、高 2 的三角形；
% 俯视图为半径 2 的圆，内含半径 1 的虚线圆（下部圆柱被遮挡的轮廓）。
% ============================================================
\begin{tikzpicture}[>=Stealth, line width=0.6pt, font=\small, scale=0.55,
  line cap=round, line join=round]
  % ---------- 正视图 ----------
  % 下部矩形：x∈[1,3]，高 4；上部三角形：底 x∈[0,4] 于 y=4，顶点 (2,6)
  \draw (1,0) -- (3,0) -- (3,4) -- (4,4) -- (2,6) -- (0,4) -- (1,4) -- (1,0);
  % 底宽 2（矩形）
  \draw[<->] (1,-0.5) -- (3,-0.5);
  \node[below] at (2,-0.5) {$2$};
  % 底宽 4（三角形底边）
  \draw[<->] (-0.6,-1.3) -- (4.6,-1.3);
  \draw (-0.6,-0.9) -- (-0.6,-1.7);
  \draw (4.6,-0.9) -- (4.6,-1.7);
  \node[below] at (2,-1.3) {$4$};
  \node[below] at (2,-2.2) {正视图};
  % 柱高 4、锥高 2（置于右侧）
  \draw[<->] (5.4,0) -- (5.4,4);
  \draw (5.0,0) -- (5.8,0);
  \draw (5.0,4) -- (5.8,4);
  \node[right] at (5.4,2) {$4$};
  \draw[<->] (5.4,4) -- (5.4,6);
  \draw (5.0,6) -- (5.8,6);
  \node[right] at (5.4,5) {$2$};

  % ---------- 侧视图 ----------
  \begin{scope}[shift={(9.2,0)}]
    \draw (1,0) -- (3,0) -- (3,4) -- (4,4) -- (2,6) -- (0,4) -- (1,4) -- (1,0);
    \draw[<->] (1,-0.5) -- (3,-0.5);
    \node[below] at (2,-0.5) {$2$};
    \draw[<->] (-0.6,-1.3) -- (4.6,-1.3);
    \draw (-0.6,-0.9) -- (-0.6,-1.7);
    \draw (4.6,-0.9) -- (4.6,-1.7);
    \node[below] at (2,-1.3) {$4$};
    \node[below] at (2,-2.2) {侧视图};
  \end{scope}

  % ---------- 俯视图 ----------
  \begin{scope}[shift={(2,-6.6)}]
    \draw (0,0) circle (2);
    \draw[dashed] (0,0) circle (1);
    \node[below] at (0,-2.7) {俯视图};
  \end{scope}
\end{tikzpicture}
\end{center}
\nopagebreak[4]
\Needspace{4\baselineskip}
\textbf{解析}

由三视图，该几何体由上、下两部分同轴叠放而成：

\begin{itemize}
\item 上部是圆锥：正视图、侧视图中的三角形底宽为 $4$、高为 $2$，故底面半径 $r_1=2$、高 $h_1=2$，体积
\end{itemize}
\nopagebreak[4]
\[V_1=\frac13\pi r_1^2h_1=\frac13\pi\cdot4\cdot2=\frac{8\pi}{3};\]
\begin{itemize}
\item 下部是圆柱：矩形宽为 $2$、高为 $4$，故底面半径 $r_2=1$、高 $h_2=4$，体积
\end{itemize}
\nopagebreak[4]
\[V_2=\pi r_2^2h_2=\pi\cdot1\cdot4=4\pi.\]

\Needspace{5\baselineskip}
俯视图中外圆半径 $2$ 对应圆锥底面，内虚线圆半径 $1$ 对应被遮挡的圆柱，与上述判断一致。故总体积为
\nopagebreak[4]
\[V=V_1+V_2=\frac{8\pi}{3}+4\pi=\frac{20\pi}{3}.\]

填 $\dfrac{20\pi}{3}$。

\gkanswer{答案：$\dfrac{20\pi}{3}$。几何体为半径 $2$、高 $2$ 的圆锥与半径 $1$、高 $4$ 的圆柱同轴叠放，$\dfrac{8\pi}{3}+4\pi=\dfrac{20\pi}{3}$。}

\end{answer}
\begin{answer}
\gksolhead{11}{}

\textbf{题目}：

设 $\{a_n\}$ 是首项为 $a_1$，公差为 $-1$ 的等差数列，$S_n$ 为其前 $n$ 项和。若 $S_1$，$S_2$，$S_4$ 成等比数列，则 $a_1$ 的值为\gkblank{}。

\Needspace{4\baselineskip}
\textbf{解析}

\Needspace{5\baselineskip}
公差 $d=-1$，故
\nopagebreak[4]
\[S_1=a_1,\qquad S_2=2a_1-1,\qquad S_4=4a_1+\frac{4\times3}{2}(-1)=4a_1-6.\]
\Needspace{5\baselineskip}
$S_1$，$S_2$，$S_4$ 成等比数列，且 $S_1=a_1\ne0$，故
\nopagebreak[4]
\[S_2^2=S_1S_4,\]
\Needspace{5\baselineskip}
即
\nopagebreak[4]
\[(2a_1-1)^2=a_1(4a_1-6).\]
\Needspace{5\baselineskip}
展开：
\nopagebreak[4]
\[4a_1^2-4a_1+1=4a_1^2-6a_1,\]
\Needspace{5\baselineskip}
所以
\nopagebreak[4]
\[2a_1=-1,\qquad a_1=-\frac12.\]
检验：此时 $S_1=-\dfrac12$，$S_2=-2$，$S_4=-8$，公比为 $4$，确成等比数列。

填 $-\dfrac12$。

\gkanswer{答案：$-\dfrac12$。由 $S_2^2=S_1S_4$ 得 $(2a_1-1)^2=a_1(4a_1-6)$，解得 $a_1=-\dfrac12$。}

\end{answer}
\begin{answer}
\gksolhead{12}{}

\textbf{题目}：

在 $\triangle ABC$ 中，内角 $A$，$B$，$C$ 所对的边分别是 $a$，$b$，$c$，已知 $b-c=\dfrac14a$，$2\sin B=3\sin C$，则 $\cos A$ 的值为\gkblank{}。

\Needspace{4\baselineskip}
\textbf{解析}

\Needspace{5\baselineskip}
由正弦定理 $\dfrac{b}{\sin B}=\dfrac{c}{\sin C}$，结合 $2\sin B=3\sin C$ 得
\nopagebreak[4]
\[2b=3c,\qquad b=\frac32c.\]
\Needspace{5\baselineskip}
代入 $b-c=\dfrac14a$：
\nopagebreak[4]
\[\frac32c-c=\frac12c=\frac14a,\]
\Needspace{5\baselineskip}
故
\nopagebreak[4]
\[a=2c.\]
\Needspace{5\baselineskip}
由余弦定理
\nopagebreak[4]
\[\begin{aligned}
\cos A&=\frac{b^2+c^2-a^2}{2bc}\\
&=\frac{\frac94c^2+c^2-4c^2}{2\cdot\frac32c\cdot c}\\
&=\frac{-\frac34c^2}{3c^2}\\
&=-\frac14.
\end{aligned}\]

填 $-\dfrac14$。

\gkanswer{答案：$-\dfrac14$。由 $2\sin B=3\sin C$ 得 $b=\dfrac32c$，再由 $b-c=\dfrac a4$ 得 $a=2c$，用余弦定理得 $\cos A=-\dfrac14$。}

\end{answer}
\begin{answer}
\gksolhead{13}{}

\textbf{题目}：

在以 $O$ 为极点的极坐标系中，圆 $\rho=4\sin\theta$ 和直线 $\rho\sin\theta=a$ 相交于 $A$，$B$ 两点。若 $\triangle AOB$ 是等边三角形，则 $a$ 的值为\gkblank{}。

\Needspace{4\baselineskip}
\textbf{解析}

\Needspace{5\baselineskip}
由 $\rho=4\sin\theta$ 得 $\rho^2=4\rho\sin\theta$，即
\nopagebreak[4]
\[x^2+y^2=4y,\qquad x^2+(y-2)^2=4.\]
\Needspace{5\baselineskip}
这是圆心 $(0,2)$、半径 $2$ 的圆。直线 $\rho\sin\theta=a$ 即 $y=a$，设交点
\nopagebreak[4]
\[A(x_0,a),\qquad B(-x_0,a),\qquad x_0>0.\]
\Needspace{5\baselineskip}
代入圆方程：
\nopagebreak[4]
\[x_0^2+(a-2)^2=4,\qquad x_0^2=4-(a-2)^2.\]
\Needspace{5\baselineskip}
于是
\nopagebreak[4]
\[|AB|=2x_0,\qquad |OA|^2=x_0^2+a^2=4-(a-2)^2+a^2=4a.\]
\Needspace{5\baselineskip}
$\triangle AOB$ 为等边三角形，故 $|AB|=|OA|$，即
\nopagebreak[4]
\[4x_0^2=4a,\qquad x_0^2=a.\]
\Needspace{5\baselineskip}
代入 $x_0^2=4-(a-2)^2=4a-a^2$：
\nopagebreak[4]
\[4a-a^2=a,\qquad a^2-3a=0,\qquad a(a-3)=0.\]
又 $a>0$（否则圆与直线不相交于两个不同点），故 $a=3$。此时 $x_0^2=3>0$，确有两个交点。

填 $3$。

\gkanswer{答案：$3$。圆为 $x^2+(y-2)^2=4$，弦 $y=a$ 的半弦长 $\sqrt{4-(a-2)^2}$；由 $|AB|=|OA|$ 得 $4-(a-2)^2=a$，即 $a=3$。}

\end{answer}
\begin{answer}
\gksolhead{14}{}

\textbf{题目}：

已知函数 $f(x)=|x^2+3x|$，$x\in\mathbf R$。若方程 $f(x)-a|x-1|=0$ 恰有 $4$ 个互异的实数根，则实数 $a$ 的取值范围为\gkblank{}。

\Needspace{4\baselineskip}
\textbf{解析}

\Needspace{5\baselineskip}
方程即
\nopagebreak[4]
\[|x^2+3x|=a|x-1|.\]
若 $a\leqslant0$，则 $a|x-1|\leqslant0$，而 $|x^2+3x|\geqslant0$，方程至多在 $|x^2+3x|=0$ 且 $a=0$ 时成立，不可能有 $4$ 个互异实根。以下设 $a>0$。
\Needspace{5\baselineskip}
由于 $x=0$ 时左边为 $0$、右边为 $a$；$x=1$ 时左边为 $4$、右边为 $0$；$x=-3$ 时左边为 $0$、右边为 $4a$，故在 $a>0$ 时 $x=0,1,-3$ 均不是根。按区间讨论，并记
\nopagebreak[4]
\[p(x)=x^2+(3+a)x-a,\qquad q(x)=x^2+(3-a)x+a.\]

\Needspace{5\baselineskip}
\textbf{（一）$x\leqslant-3$。} 此时 $x^2+3x\geqslant0$，$|x-1|=1-x$，方程化为
\nopagebreak[4]
\[x^2+3x=a(1-x)\iff p(x)=0.\]
$p(-3)=9-3(3+a)-a=-4a<0$，而当 $x\to-\infty$ 时 $p(x)\to+\infty$；又 $p$ 的常数项为 $-a<0$，两根异号，故 $p$ 在 $(-\infty,-3)$ 内恰有一个根。此区间贡献 $1$ 个根。

\textbf{（二）$0\leqslant x<1$。} 此时 $x^2+3x\geqslant0$，$|x-1|=1-x$，方程仍为 $p(x)=0$。$p(0)=-a<0$，$p(1)=4>0$，故在 $(0,1)$ 内恰有一个根（另一根为负）。此区间贡献 $1$ 个根。

\Needspace{5\baselineskip}
\textbf{（三）$-3<x<0$。} 此时 $x^2+3x<0$，$|x-1|=1-x$，方程化为
\nopagebreak[4]
\[-(x^2+3x)=a(1-x)\iff q(x)=0.\]
\Needspace{5\baselineskip}
其判别式
\nopagebreak[4]
\[\Delta=(3-a)^2-4a=a^2-10a+9=(a-1)(a-9).\]
又 $q(-3)=9-3(3-a)+a=4a>0$，$q(0)=a>0$，顶点横坐标为 $\dfrac{a-3}{2}$，两根之和为 $a-3$、两根之积为 $a$。

\begin{itemize}
\item 当 $0<a<1$ 时，$\Delta>0$，和 $a-3\in(-3,-2)$ 为负、积 $a>0$，故两根均为负；结合 $q(-3)>0$、$q(0)>0$ 与顶点处 $q<0$，顶点 $\dfrac{a-3}{2}\in\left(-1,-\dfrac12\right)$ 落在区间内，所以两根都在 $(-3,0)$ 内，贡献 $2$ 个根。
\item 当 $a=1$ 时，$q(x)=x^2+2x+1=(x+1)^2$，在 $(-3,0)$ 内有重根 $x=-1$，贡献 $1$ 个根。
\item 当 $1<a<9$ 时，$\Delta<0$，无实根，贡献 $0$ 个。
\item 当 $a=9$ 时，$q(x)=x^2-6x+9=(x-3)^2$，根 $x=3\notin(-3,0)$，贡献 $0$ 个。
\item 当 $a>9$ 时，$\Delta>0$，和 $a-3>6>0$、积 $a>0$，两根均为正，不在 $(-3,0)$ 内，贡献 $0$ 个。
\end{itemize}

\Needspace{5\baselineskip}
\textbf{（四）$x>1$。} 此时 $x^2+3x>0$，$|x-1|=x-1$，方程化为
\nopagebreak[4]
\[x^2+3x=a(x-1)\iff q(x)=0.\]
注意 $q(1)=1+(3-a)+a=4>0$。

\begin{itemize}
\item 当 $a>9$ 时，$\Delta>0$，顶点 $\dfrac{a-3}{2}>3>1$，和 $a-3>6>0$、积 $a>0$ 说明两根同为正，再由 $q(1)>0$ 知两根都大于 $1$，贡献 $2$ 个根。
\item 当 $a=9$ 时，$q(x)=(x-3)^2$，根 $x=3>1$，贡献 $1$ 个根。
\item 当 $a<9$ 时：若 $1<a<9$ 则 $\Delta<0$ 无根；若 $a=1$ 则重根 $-1\not>1$；若 $0<a<1$ 则 $\Delta>0$，但两根之和 $a-3<0$、两根之积 $a>0$，两根均为负，更不可能大于 $1$。总之贡献 $0$ 个根。
\end{itemize}

各区间根数汇总：

\begin{center}
\begin{tabular}{llllll}
\toprule
$a$ 的范围 & $x\leqslant-3$ & $0\leqslant x<1$ & $-3<x<0$ & $x>1$ & 合计 \\
\midrule
$0<a<1$ & $1$ & $1$ & $2$ & $0$ & $4$ \\
$a=1$ & $1$ & $1$ & $1$ & $0$ & $3$ \\
$1<a<9$ & $1$ & $1$ & $0$ & $0$ & $2$ \\
$a=9$ & $1$ & $1$ & $0$ & $1$ & $3$ \\
$a>9$ & $1$ & $1$ & $0$ & $2$ & $4$ \\
\bottomrule
\end{tabular}
\end{center}

\Needspace{5\baselineskip}
恰有 $4$ 个互异实根当且仅当
\nopagebreak[4]
\[a\in(0,1)\cup(9,+\infty).\]

填 $(0,1)\cup(9,+\infty)$。

\gkanswer{答案：$(0,1)\cup(9,+\infty)$。在 $(-3,0)$ 内根的个数为 $0<a<1$ 时 $2$ 个、$a=1$ 时 $1$ 个、$a>1$ 时 $0$ 个；在 $(1,+\infty)$ 内为 $a>9$ 时 $2$ 个、$a=9$ 时 $1$ 个、$a<9$ 时 $0$ 个；另两支各恒有 $1$ 个根，合计恰为 $4$ 个当且仅当 $a\in(0,1)\cup(9,+\infty)$。}

\end{answer}
\begin{answer}
\gksolhead{15}{}

\textbf{题目}：

已知函数 $f(x)=\cos x\cdot\sin\left(x+\dfrac\pi3\right)-\sqrt3\cos^2x+\dfrac{\sqrt3}{4}$，$x\in\mathbf R$。

（1） 求 $f(x)$ 的最小正周期；

（2） 求 $f(x)$ 在闭区间 $\left[-\dfrac\pi4,\dfrac\pi4\right]$ 上的最大值和最小值。

\Needspace{4\baselineskip}
\textbf{解析}

\Needspace{5\baselineskip}
先用和角公式展开：
\nopagebreak[4]
\[\begin{aligned}
f(x)&=\cos x\left(\sin x\cos\frac\pi3+\cos x\sin\frac\pi3\right)-\sqrt3\cos^2x+\frac{\sqrt3}{4}\\
&=\cos x\left(\frac12\sin x+\frac{\sqrt3}{2}\cos x\right)-\sqrt3\cos^2x+\frac{\sqrt3}{4}\\
&=\frac12\sin x\cos x+\frac{\sqrt3}{2}\cos^2x-\sqrt3\cos^2x+\frac{\sqrt3}{4}\\
&=\frac14\sin2x-\frac{\sqrt3}{2}\cos^2x+\frac{\sqrt3}{4}\\
&=\frac14\sin2x-\frac{\sqrt3}{4}(1+\cos2x)+\frac{\sqrt3}{4}\\
&=\frac14\sin2x-\frac{\sqrt3}{4}\cos2x\\
&=\frac12\left(\frac12\sin2x-\frac{\sqrt3}{2}\cos2x\right)\\
&=\frac12\sin\left(2x-\frac\pi3\right).
\end{aligned}\]

\gksubpart{（1）}

\Needspace{5\baselineskip}
由 $f(x)=\dfrac12\sin\left(2x-\dfrac\pi3\right)$，最小正周期为
\nopagebreak[4]
\[T=\frac{2\pi}{2}=\pi.\]

\gksubpart{（2）}

\Needspace{5\baselineskip}
当 $x\in\left[-\dfrac\pi4,\dfrac\pi4\right]$ 时，
\nopagebreak[4]
\[2x-\frac\pi3\in\left[-\frac{5\pi}{6},\frac\pi6\right].\]
在此区间上 $\sin$ 的最小值为 $-1$，在 $2x-\dfrac\pi3=-\dfrac\pi2$（即 $x=-\dfrac{\pi}{12}$）处取到；最大值为 $\dfrac12$，在区间右端点 $2x-\dfrac\pi3=\dfrac\pi6$（即 $x=\dfrac\pi4$）处取到。

\Needspace{5\baselineskip}
故
\nopagebreak[4]
\[f(x)_{\max}=\frac12\cdot\frac12=\frac14\quad\left(x=\frac\pi4\right),\]
\nopagebreak[4]
\[f(x)_{\min}=\frac12\cdot(-1)=-\frac12\quad\left(x=-\frac{\pi}{12}\right).\]

\gkanswer{答案（1）：$\pi$。化简得 $f(x)=\dfrac12\sin\left(2x-\dfrac\pi3\right)$，故最小正周期为 $\dfrac{2\pi}{2}=\pi$。}

\gkanswer{答案（2）：最大值 $\dfrac14$，最小值 $-\dfrac12$。当 $2x-\dfrac\pi3\in\left[-\dfrac{5\pi}{6},\dfrac\pi6\right]$ 时，$\sin$ 的最大值为 $\dfrac12$（$x=\dfrac\pi4$），最小值为 $-1$（$x=-\dfrac{\pi}{12}$）。}

\end{answer}
\begin{answer}
\gksolhead{16}{}

\textbf{题目}：

某大学志愿者协会有 $6$ 名男同学，$4$ 名女同学。在这 $10$ 名同学中，$3$ 名同学来自数学学院，其余 $7$ 名同学来自物理、化学等其他互不相同的七个学院。现从这 $10$ 名同学中随机选取 $3$ 名同学，到希望小学进行支教活动（每位同学被选到的可能性相同）。

（1） 求选出的 $3$ 名同学是来自互不相同学院的概率；

（2） 设 $X$ 为选出的 $3$ 名同学中女同学的人数，求随机变量 $X$ 的分布列和数学期望。

\Needspace{4\baselineskip}
\textbf{解析}

\Needspace{5\baselineskip}
从 $10$ 人中任取 $3$ 人，总数为
\nopagebreak[4]
\[\mathrm C_{10}^3=\frac{10\times9\times8}{3\times2\times1}=120,\]
且每种选法等可能。

\gksubpart{（1）}

$3$ 名数学学院同学同属一个学院，其余 $7$ 人分属七个互不相同的学院。要使选出的 $3$ 人来自互不相同的学院：

\begin{itemize}
\item 不选数学学院的人：从其余 $7$ 个学院各至多取 $1$ 人，有 $\mathrm C_7^3=35$ 种；
\item 选 $1$ 名数学学院的人：$\mathrm C_3^1$ 种，再从其余 $7$ 个学院中选 $2$ 个不同学院，有 $\mathrm C_7^2$ 种，共
\end{itemize}
\nopagebreak[4]
\[3\times21=63\]
种。

\Needspace{5\baselineskip}
选 $2$ 名或 $3$ 名数学学院的人时，至少两人同属数学学院，不符合要求。故有利结果数为
\nopagebreak[4]
\[35+63=98,\]
\Needspace{5\baselineskip}
所求概率为
\nopagebreak[4]
\[P=\frac{98}{120}=\frac{49}{60}.\]

\gksubpart{（2）}

\Needspace{5\baselineskip}
$X$ 的取值为 $0,1,2,3$，按超几何分布
\nopagebreak[4]
\[P(X=k)=\frac{\mathrm C_4^k\mathrm C_6^{3-k}}{\mathrm C_{10}^3},\qquad k=0,1,2,3.\]
\Needspace{5\baselineskip}
逐项计算：
\nopagebreak[4]
\[P(X=0)=\frac{\mathrm C_4^0\mathrm C_6^3}{120}=\frac{20}{120}=\frac16,\]
\nopagebreak[4]
\[P(X=1)=\frac{\mathrm C_4^1\mathrm C_6^2}{120}=\frac{4\times15}{120}=\frac{60}{120}=\frac12,\]
\nopagebreak[4]
\[P(X=2)=\frac{\mathrm C_4^2\mathrm C_6^1}{120}=\frac{6\times6}{120}=\frac{36}{120}=\frac3{10},\]
\nopagebreak[4]
\[P(X=3)=\frac{\mathrm C_4^3\mathrm C_6^0}{120}=\frac{4}{120}=\frac1{30}.\]
检验：$\dfrac16+\dfrac12+\dfrac3{10}+\dfrac1{30}=\dfrac{20+60+36+4}{120}=1$。

分布列为

\begin{center}
\begin{tabular}{lllll}
\toprule
$X$ & $0$ & $1$ & $2$ & $3$ \\
\midrule
$P$ & $\dfrac16$ & $\dfrac12$ & $\dfrac3{10}$ & $\dfrac1{30}$ \\
\bottomrule
\end{tabular}
\end{center}

\Needspace{5\baselineskip}
数学期望
\nopagebreak[4]
\[\begin{aligned}
E(X)&=0\times\frac16+1\times\frac12+2\times\frac3{10}+3\times\frac1{30}\\
&=\frac12+\frac35+\frac1{10}\\
&=\frac{5+6+1}{10}=\frac65.
\end{aligned}\]
也可由超几何分布的期望公式 $E(X)=3\times\dfrac4{10}=\dfrac65$ 直接得到。

\gkanswer{答案（1）：$\dfrac{49}{60}$。数学学院的 $3$ 人同属一个学院，其余 $7$ 人各属不同学院；符合条件的结果数为 $\mathrm C_7^3+3\mathrm C_7^2=98$，而总数为 $\mathrm C_{10}^3=120$。}

\gkanswer{答案（2）：$P(X=0)=\dfrac16$，$P(X=1)=\dfrac12$，$P(X=2)=\dfrac3{10}$，$P(X=3)=\dfrac1{30}$，$E(X)=\dfrac65$。按超几何分布计算，或直接取 $E(X)=3\times\dfrac4{10}=\dfrac65$。}

\end{answer}
\begin{answer}
\gksolhead{17}{}

\textbf{题目}：

如图，在四棱锥 $P\text{-}ABCD$ 中，$PA\perp$ 底面 $ABCD$，$AD\perp AB$，$AB\parallel DC$，$AD=DC=AP=2$，$AB=1$，点 $E$ 为棱 $PC$ 的中点。
\nopagebreak[4]
\begin{center}% ============================================================
% f2014_tianjin_science_q17_pyramid.tex
% 2014 天津卷（理科）第 17 题：四棱锥 P-ABCD
% 来源 PDF：普通高考/2014/2014天津理.pdf
% 原页：第 2 页（A3），印面 924，第 17 题题干下方
% PA⊥底面 ABCD，AD⊥AB，AB∥DC，AD=DC=AP=2，AB=1，E 为 PC 中点。
% 由三视图式平行投影绘制，不按严格比例；PB、BE 为可见棱，
% PD、DC、AD、AC、BD 为不可见线（虚线）。
% ============================================================
\begin{tikzpicture}[>=Stealth, line width=0.65pt, font=\small, scale=1.0,
  line cap=round, line join=round]
  \coordinate (A) at (0,0);
  \coordinate (B) at (2.00,0);
  \coordinate (D) at (1.76,1.47);
  \coordinate (C) at (5.76,1.47);
  \coordinate (P) at (0,3.94);
  \coordinate (E) at (2.88,2.705);

  % 可见棱
  \draw (P) -- (A);
  \draw (P) -- (B);
  \draw (P) -- (C);
  \draw (A) -- (B);
  \draw (B) -- (C);
  \draw (B) -- (E);

  % 不可见线
  \draw[dashed] (P) -- (D);
  \draw[dashed] (A) -- (D);
  \draw[dashed] (D) -- (C);
  \draw[dashed] (A) -- (C);
  \draw[dashed] (B) -- (D);

  \fill (A) circle (0.85pt) node[below left=-1pt] {$A$};
  \fill (B) circle (0.85pt) node[below=-1pt] {$B$};
  \fill (C) circle (0.85pt) node[right=-1pt] {$C$};
  \fill (D) circle (0.85pt) node[above left=-1pt] {$D$};
  \fill (P) circle (0.85pt) node[above=-1pt] {$P$};
  \fill (E) circle (0.85pt) node[above right=-2pt] {$E$};
\end{tikzpicture}
\end{center}
\nopagebreak[4]
（1） 证明：$BE\perp DC$；

（2） 求直线 $BE$ 与平面 $PBD$ 所成角的正弦值；

（3） 若 $F$ 为棱 $PC$ 上一点，满足 $BF\perp AC$，求二面角 $F\text{-}AB\text{-}P$ 的余弦值。

\Needspace{4\baselineskip}
\textbf{解析}

以 $A$ 为原点，$\overrightarrow{AB}$ 的方向为 $x$ 轴正方向，$\overrightarrow{AD}$ 的方向为 $y$ 轴正方向，$\overrightarrow{AP}$ 的方向为 $z$ 轴正方向，建立空间直角坐标系。

\Needspace{5\baselineskip}
由 $AB=1$，$AD=2$，$AP=2$ 得
\nopagebreak[4]
\[A(0,0,0),\quad B(1,0,0),\quad D(0,2,0),\quad P(0,0,2).\]
\Needspace{5\baselineskip}
又 $DC\parallel AB$ 且 $DC=2$，故
\nopagebreak[4]
\[C=(2,2,0).\]
\Needspace{5\baselineskip}
$E$ 为 $PC$ 的中点：
\nopagebreak[4]
\[E\left(\frac{2+0}{2},\frac{2+0}{2},\frac{0+2}{2}\right)=(1,1,1).\]

\gksubpart{（1）}

\nopagebreak[4]
\[\overrightarrow{BE}=(1-1,1-0,1-0)=(0,1,1),\]
\nopagebreak[4]
\[\overrightarrow{DC}=(2-0,2-2,0-0)=(2,0,0).\]
\Needspace{5\baselineskip}
于是
\nopagebreak[4]
\[\overrightarrow{BE}\cdot\overrightarrow{DC}=0\times2+1\times0+1\times0=0,\]
所以 $BE\perp DC$。

\gksubpart{（2）}

\nopagebreak[4]
\[\overrightarrow{PB}=(1,0,-2),\qquad \overrightarrow{PD}=(0,2,-2).\]
\Needspace{5\baselineskip}
取平面 $PBD$ 的法向量
\nopagebreak[4]
\[\boldsymbol n=\overrightarrow{PB}\times\overrightarrow{PD}
=(0\cdot(-2)-(-2)\cdot2,\ (-2)\cdot0-1\cdot(-2),\ 1\cdot2-0\cdot0)
=(4,2,2),\]
可化简为 $(2,1,1)$，且 $|\boldsymbol n|=\sqrt{2^2+1^2+1^2}=\sqrt6$。

\Needspace{5\baselineskip}
又 $|\overrightarrow{BE}|=\sqrt2$，
\nopagebreak[4]
\[\overrightarrow{BE}\cdot\boldsymbol n=0\times2+1\times1+1\times1=2.\]
\Needspace{5\baselineskip}
设直线 $BE$ 与平面 $PBD$ 所成角为 $\theta$，则
\nopagebreak[4]
\[\sin\theta=\frac{|\overrightarrow{BE}\cdot\boldsymbol n|}{|\overrightarrow{BE}|\,|\boldsymbol n|}
=\frac{2}{\sqrt2\cdot\sqrt6}=\frac{2}{2\sqrt3}=\frac{\sqrt3}{3}.\]

\gksubpart{（3）}

\Needspace{5\baselineskip}
设 $F$ 在棱 $PC$ 上，且
\nopagebreak[4]
\[\overrightarrow{PF}=t\,\overrightarrow{PC},\qquad 0\leqslant t\leqslant1.\]
\Needspace{5\baselineskip}
由 $\overrightarrow{PC}=(2,2,-2)$ 得
\nopagebreak[4]
\[F=(2t,2t,2-2t),\]
\Needspace{5\baselineskip}
从而
\nopagebreak[4]
\[\overrightarrow{BF}=(2t-1,2t,2-2t).\]
\Needspace{5\baselineskip}
又 $\overrightarrow{AC}=(2,2,0)$。由 $BF\perp AC$ 得
\nopagebreak[4]
\[\overrightarrow{BF}\cdot\overrightarrow{AC}=2(2t-1)+2\cdot2t+0=8t-2=0,\]
\Needspace{5\baselineskip}
故 $t=\dfrac14$，
\nopagebreak[4]
\[F\left(\frac12,\frac12,\frac32\right).\]

\Needspace{5\baselineskip}
平面 $ABP$ 即 $y=0$ 平面，取其法向量
\nopagebreak[4]
\[\boldsymbol n_1=(0,1,0).\]
\Needspace{5\baselineskip}
平面 $ABF$ 中，$\overrightarrow{AB}=(1,0,0)$，$\overrightarrow{AF}=\left(\dfrac12,\dfrac12,\dfrac32\right)$，取法向量
\nopagebreak[4]
\[\boldsymbol n_2=\overrightarrow{AB}\times\overrightarrow{AF}
=\left(0\cdot\frac32-0\cdot\frac12,\ 0\cdot\frac12-1\cdot\frac32,\ 1\cdot\frac12-0\cdot\frac12\right)
=\left(0,-\frac32,\frac12\right),\]
可化简为 $(0,-3,1)$，且 $|\boldsymbol n_2|=\sqrt{10}$。

\Needspace{5\baselineskip}
二面角 $F\text{-}AB\text{-}P$ 的平面角在过 $A$ 且垂直于 $AB$ 的平面 $x=0$ 内：半平面 $ABF$ 截得射线方向可取 $(0,1,3)$，半平面 $ABP$ 截得射线方向可取 $(0,0,1)$，两方向夹角为锐角，其余弦为
\nopagebreak[4]
\[\cos\theta=\frac{(0,1,3)\cdot(0,0,1)}{\sqrt{10}\cdot1}=\frac{3}{\sqrt{10}}=\frac{3\sqrt{10}}{10}.\]
（用两个法向量计算时 $|\boldsymbol n_1\cdot\boldsymbol n_2|=3$，同样得到余弦的绝对值为 $\dfrac{3}{\sqrt{10}}$。）

所以二面角 $F\text{-}AB\text{-}P$ 的余弦值为 $\dfrac{3\sqrt{10}}{10}$。

\gkanswer{答案（1）：见解析版}

\gkanswer{答案（2）：$\dfrac{\sqrt3}{3}$。平面 $PBD$ 的法向量取 $(2,1,1)$，$\sin\theta=\dfrac{|\overrightarrow{BE}\cdot\boldsymbol n|}{|\overrightarrow{BE}||\boldsymbol n|}=\dfrac{4}{\sqrt2\cdot2\sqrt6}=\dfrac{\sqrt3}{3}$。}

\gkanswer{答案（3）：$\dfrac{3\sqrt{10}}{10}$。由 $\overrightarrow{BF}\cdot\overrightarrow{AC}=0$ 得 $F$ 分 $PC$ 的比为 $\dfrac14$，两半平面法向量分别取 $(0,1,0)$ 与 $(0,-3,1)$，其余弦绝对值为 $\dfrac{3}{\sqrt{10}}$，二面角为锐角。}

\end{answer}
\begin{answer}
\gksolhead{18}{}

\textbf{题目}：

设椭圆 $\dfrac{x^2}{a^2}+\dfrac{y^2}{b^2}=1$（$a>b>0$）的左、右焦点分别为 $F_1$，$F_2$，右顶点为 $A$，上顶点为 $B$，已知 $|AB|=\dfrac{\sqrt3}{2}|F_1F_2|$。

（1） 求椭圆的离心率；

（2） 设 $P$ 为椭圆上异于其顶点的一点，以线段 $PB$ 为直径的圆经过点 $F_1$，经过原点 $O$ 的直线 $l$ 与该圆相切，求直线 $l$ 的斜率。

\Needspace{4\baselineskip}
\textbf{解析}

\gksubpart{（1）}

\Needspace{5\baselineskip}
$A(a,0)$，$B(0,b)$，故
\nopagebreak[4]
\[|AB|^2=a^2+b^2,\]
\Needspace{5\baselineskip}
而 $|F_1F_2|=2c$。由题设
\nopagebreak[4]
\[a^2+b^2=\left(\frac{\sqrt3}{2}\right)^2(2c)^2=3c^2.\]
\Needspace{5\baselineskip}
又 $b^2=a^2-c^2$，代入得
\nopagebreak[4]
\[2a^2-c^2=3c^2,\qquad a^2=2c^2.\]
\Needspace{5\baselineskip}
所以离心率
\nopagebreak[4]
\[e=\frac ca=\frac1{\sqrt2}=\frac{\sqrt2}{2}.\]

\gksubpart{（2）}

\Needspace{5\baselineskip}
由（1）$a^2=2c^2$，故 $b^2=a^2-c^2=c^2$，椭圆为
\nopagebreak[4]
\[\frac{x^2}{2c^2}+\frac{y^2}{c^2}=1,\]
\Needspace{5\baselineskip}
即
\nopagebreak[4]
\[x^2+2y^2=2c^2.\]
此时 $B(0,c)$，$F_1(-c,0)$。

\Needspace{5\baselineskip}
因为以 $PB$ 为直径的圆经过 $F_1$，所以 $\angle PF_1B=90^\circ$，即
\nopagebreak[4]
\[\overrightarrow{F_1P}\cdot\overrightarrow{F_1B}=0.\]
\Needspace{5\baselineskip}
设 $P(x,y)$，则 $\overrightarrow{F_1P}=(x+c,y)$，$\overrightarrow{F_1B}=(c,c)$，故
\nopagebreak[4]
\[c(x+c)+cy=0,\qquad x+y=-c. \tag{1}\]
\Needspace{5\baselineskip}
联立 $x^2+2y^2=2c^2$ 与（1）：由 $x=-c-y$ 得
\nopagebreak[4]
\[(y+c)^2+2y^2=2c^2,\qquad 3y^2+2cy-c^2=0.\]
\Needspace{5\baselineskip}
解得
\nopagebreak[4]
\[y=\frac{-2c\pm\sqrt{4c^2+12c^2}}{6}=\frac{-2c\pm4c}{6},\]
即 $y=\dfrac c3$ 或 $y=-c$。

\Needspace{5\baselineskip}
当 $y=-c$ 时 $x=0$，$P(0,-c)$ 是椭圆的下顶点，与“$P$ 异于其顶点”矛盾。故
\nopagebreak[4]
\[y=\frac c3,\qquad x=-c-\frac c3=-\frac{4c}{3},\]
\Needspace{5\baselineskip}
即
\nopagebreak[4]
\[P\left(-\frac{4c}{3},\frac c3\right).\]

\Needspace{5\baselineskip}
以 $PB$ 为直径的圆，圆心 $C'$ 为 $PB$ 的中点：
\nopagebreak[4]
\[C'\left(-\frac{2c}{3},\frac{2c}{3}\right).\]
\Needspace{5\baselineskip}
半径
\nopagebreak[4]
\[r=\frac12|PB|=\frac12\sqrt{\left(-\frac{4c}{3}\right)^2+\left(\frac c3-c\right)^2}
=\frac12\sqrt{\frac{16c^2}{9}+\frac{4c^2}{9}}
=\frac12\cdot\frac{2\sqrt5c}{3}=\frac{\sqrt5c}{3}.\]

\Needspace{5\baselineskip}
直线 $l$ 经过原点，若 $l$ 为 $x=0$，则圆心到它的距离为 $\dfrac{2c}{3}\ne r$，不能相切；故可设
\nopagebreak[4]
\[l:\ y=kx,\qquad kx-y=0.\]
\Needspace{5\baselineskip}
相切时圆心到 $l$ 的距离等于半径：
\nopagebreak[4]
\[\frac{\left|k\left(-\frac{2c}{3}\right)-\frac{2c}{3}\right|}{\sqrt{k^2+1}}=\frac{\sqrt5c}{3},\]
\Needspace{5\baselineskip}
即
\nopagebreak[4]
\[\frac{2c}{3}\cdot\frac{|k+1|}{\sqrt{k^2+1}}=\frac{\sqrt5c}{3},\]
\nopagebreak[4]
\[2|k+1|=\sqrt5\sqrt{k^2+1}.\]
\Needspace{5\baselineskip}
两边平方：
\nopagebreak[4]
\[4(k+1)^2=5(k^2+1),\]
\nopagebreak[4]
\[4k^2+8k+4=5k^2+5,\qquad k^2-8k+1=0.\]
\Needspace{5\baselineskip}
解得
\nopagebreak[4]
\[k=\frac{8\pm\sqrt{64-4}}{2}=4\pm\sqrt{15}.\]

所以直线 $l$ 的斜率为 $4+\sqrt{15}$ 或 $4-\sqrt{15}$。

\gkanswer{答案（1）：$\dfrac{\sqrt2}{2}$。由 $a^2+b^2=3c^2$ 与 $b^2=a^2-c^2$ 得 $a^2=2c^2$，故 $e=\dfrac{\sqrt2}{2}$。}

\gkanswer{答案（2）：$4+\sqrt{15}$ 或 $4-\sqrt{15}$。由 $\overrightarrow{F_1P}\cdot\overrightarrow{F_1B}=0$ 得 $P\left(-\dfrac{4c}{3},\dfrac c3\right)$，圆心为 $\left(-\dfrac{2c}{3},\dfrac{2c}{3}\right)$、半径为 $\dfrac{\sqrt5c}{3}$，令 $l:y=kx$ 用点到直线距离等于半径，得 $k^2-8k+1=0$。}

\end{answer}
\begin{answer}
\gksolhead{19}{}

\textbf{题目}：

已知 $q$ 和 $n$ 均为给定的大于 $1$ 的自然数。设集合 $M=\{0,1,2,\cdots,q-1\}$，集合 $A=\{x\mid x=x_1+x_2q+\cdots+x_nq^{n-1},\ x_i\in M,\ i=1,2,\cdots,n\}$。

（1） 当 $q=2$，$n=3$ 时，用列举法表示集合 $A$；

（2） 设 $s$，$t\in A$，$s=a_1+a_2q+\cdots+a_nq^{n-1}$，$t=b_1+b_2q+\cdots+b_nq^{n-1}$，其中 $a_i$，$b_i\in M$，$i=1,2,\cdots,n$。证明：若 $a_n<b_n$，则 $s<t$。

\Needspace{4\baselineskip}
\textbf{解析}

\gksubpart{（1）}

\Needspace{5\baselineskip}
当 $q=2$，$n=3$ 时，$M=\{0,1\}$，且
\nopagebreak[4]
\[x=x_1+2x_2+4x_3,\qquad x_1,x_2,x_3\in\{0,1\}.\]
\Needspace{5\baselineskip}
这恰是三位二进制数的展开，八个三位二进制数
\nopagebreak[4]
\[000,001,010,011,100,101,110,111\]
\Needspace{5\baselineskip}
分别对应
\nopagebreak[4]
\[0,1,2,3,4,5,6,7.\]
\Needspace{5\baselineskip}
故
\nopagebreak[4]
\[A=\{0,1,2,3,4,5,6,7\}.\]

\gksubpart{（2）}

\Needspace{5\baselineskip}
由 $s$，$t$ 的表达式，
\nopagebreak[4]
\[\begin{aligned}
t-s&=\sum_{i=1}^{n}(b_i-a_i)q^{i-1}\\
&=(b_n-a_n)q^{n-1}+\sum_{i=1}^{n-1}(b_i-a_i)q^{i-1}.
\end{aligned}\]
\Needspace{5\baselineskip}
因 $a_n$，$b_n\in M=\{0,1,\cdots,q-1\}$ 且 $a_n<b_n$，故 $b_n-a_n\geqslant1$。又对每个 $i$，
\nopagebreak[4]
\[b_i-a_i\geqslant-(q-1),\]
\Needspace{5\baselineskip}
所以
\nopagebreak[4]
\[\begin{aligned}
t-s&\geqslant q^{n-1}-(q-1)\sum_{i=1}^{n-1}q^{i-1}\\
&=q^{n-1}-(q-1)\cdot\frac{q^{n-1}-1}{q-1}\\
&=q^{n-1}-(q^{n-1}-1)\\
&=1>0.
\end{aligned}\]
故 $s<t$。

\gkanswer{答案（1）：$A=\{0,1,2,3,4,5,6,7\}$。$q=2$、$n=3$ 时 $x=x_1+2x_2+4x_3$，三系数各取 $0$ 或 $1$，共 $2^3=8$ 个互异值。}

\gkanswer{答案（2）：见解析版}

\end{answer}
\begin{answer}
\gksolhead{20}{}

\textbf{题目}：

已知函数 $f(x)=x-a\mathrm e^x$（$a\in\mathbf R$），$x\in\mathbf R$，已知函数 $y=f(x)$ 有两个零点 $x_1$，$x_2$，且 $x_1<x_2$。

（1） 求 $a$ 的取值范围；

（2） 证明 $\dfrac{x_2}{x_1}$ 随着 $a$ 的减小而增大；

（3） 证明 $x_1+x_2$ 随着 $a$ 的减小而增大。

\Needspace{4\baselineskip}
\textbf{解析}

\gksubpart{（1）}

\nopagebreak[4]
\[f'(x)=1-a\mathrm e^x.\]

若 $a\leqslant0$，则 $f'(x)=1-a\mathrm e^x>0$ 对一切 $x$ 成立，$f$ 在 $\mathbf R$ 上严格递增，至多有一个零点，与题设两个零点矛盾。

\Needspace{5\baselineskip}
故 $a>0$。此时令 $f'(x)=0$ 得
\nopagebreak[4]
\[x=-\ln a.\]
\Needspace{5\baselineskip}
当 $x<-\ln a$ 时 $f'(x)>0$，$f$ 递增；当 $x>-\ln a$ 时 $f'(x)<0$，$f$ 递减。故最大值为
\nopagebreak[4]
\[f(-\ln a)=-\ln a-a\cdot\frac1a=-\ln a-1.\]
\Needspace{5\baselineskip}
又
\nopagebreak[4]
\[\lim_{x\to-\infty}f(x)=-\infty,\qquad \lim_{x\to+\infty}f(x)=-\infty.\]
\Needspace{5\baselineskip}
$f$ 有两个零点当且仅当最大值大于 $0$：
\nopagebreak[4]
\[-\ln a-1>0,\qquad \ln a<-1,\qquad 0<a<\frac1{\mathrm e}.\]
此时在 $(-\infty,-\ln a)$ 与 $(-\ln a,+\infty)$ 内各恰有一个零点，共两个。

所以 $a$ 的取值范围是 $\left(0,\dfrac1{\mathrm e}\right)$。

\gksubpart{（2）}

\Needspace{5\baselineskip}
设
\nopagebreak[4]
\[h(x)=x\mathrm e^{-x}.\]
\Needspace{5\baselineskip}
$f(x)=0$ 即 $a=x\mathrm e^{-x}=h(x)$。求导：
\nopagebreak[4]
\[h'(x)=\mathrm e^{-x}(1-x).\]
\Needspace{5\baselineskip}
故 $h$ 在 $(-\infty,1]$ 上递增，在 $[1,+\infty)$ 上递减，且
\nopagebreak[4]
\[h(1)=\frac1{\mathrm e},\qquad \lim_{x\to+\infty}h(x)=0.\]
\Needspace{5\baselineskip}
由（1），当 $0<a<\dfrac1{\mathrm e}$ 时方程 $h(x)=a$ 的两个根满足
\nopagebreak[4]
\[0<x_1<1<x_2.\]

\Needspace{5\baselineskip}
任取
\nopagebreak[4]
\[0<a_2<a_1<\frac1{\mathrm e},\]
设对应的零点分别为 $x_1(a_1)<1<x_2(a_1)$ 与 $x_1(a_2)<1<x_2(a_2)$。

\Needspace{5\baselineskip}
在 $(0,1)$ 上 $h$ 严格递增，而
\nopagebreak[4]
\[h(x_1(a_2))=a_2<a_1=h(x_1(a_1)),\]
\Needspace{5\baselineskip}
故
\nopagebreak[4]
\[x_1(a_2)<x_1(a_1).\]
\Needspace{5\baselineskip}
在 $(1,+\infty)$ 上 $h$ 严格递减，而
\nopagebreak[4]
\[h(x_2(a_2))=a_2<a_1=h(x_2(a_1)),\]
\Needspace{5\baselineskip}
故
\nopagebreak[4]
\[x_2(a_2)>x_2(a_1).\]
\Needspace{5\baselineskip}
于是
\nopagebreak[4]
\[\frac{x_2(a_2)}{x_1(a_2)}>\frac{x_2(a_1)}{x_1(a_1)},\]
即 $a$ 减小时 $\dfrac{x_2}{x_1}$ 增大。

\gksubpart{（3）}

\Needspace{5\baselineskip}
沿用（2）的记号。对 $0<a_2<a_1<\dfrac1{\mathrm e}$，已证
\nopagebreak[4]
\[x_1(a_2)<x_1(a_1),\qquad x_2(a_2)>x_2(a_1).\]
\Needspace{5\baselineskip}
两式相加得
\nopagebreak[4]
\[x_1(a_2)+x_2(a_2)>x_1(a_1)+x_2(a_1),\]
即 $a$ 减小时 $x_1+x_2$ 增大。

\gkanswer{答案（1）：$\left(0,\dfrac1{\mathrm e}\right)$。$f'(x)=1-a\mathrm e^x$，$a>0$ 时最大值 $f(-\ln a)=-\ln a-1$，有两个零点须 $-\ln a-1>0$。}

\gkanswer{答案（2）：见解析版}

\gkanswer{答案（3）：见解析版}

\end{answer}